\documentclass{amsart}
\usepackage{amsmath,amssymb,amsthm,hyperref}
\usepackage{algorithm}
\usepackage{algpseudocode}

\theoremstyle{plain}
\newtheorem{theorem}{Theorem}
\newtheorem{lemma}{Lemma}
\newtheorem{proposition}{Proposition}
\newtheorem{corollary}{Corollary}
\theoremstyle{definition}
\newtheorem{definition}{Definition}
\newtheorem{remark}{Remark}

\newcommand{\R}{\mathbb{R}}
\newcommand{\maj}{\prec}
\newcommand{\trump}{\prec_T}

\begin{document}

\title{A finite necessary-and-sufficient criterion for catalytic majorization}
\maketitle

\section{Statement of the criterion}

Throughout, $p=(p_1,\dots,p_n)$ and $q=(q_1,\dots,q_n)$ are probability vectors
($p_i,q_i\ge 0$, $\sum_i p_i=\sum_i q_i=1$) written in non-increasing order
$p_1\ge\cdots\ge p_n$, $q_1\ge\cdots\ge q_n$. For a probability vector $v$ of any
length and a real number $t$ we use the \emph{power sums}
\[
  P_t(v)\;=\;\sum_{i:\,v_i>0} v_i^{\,t},
\]
where the sum is over the strictly positive entries (so $P_t(v)$ is finite for every
real $t$, including $t<0$, and $P_0(v)=d(v):=\#\{i:v_i>0\}$, $P_1(v)=1$). We write
\[
  H(v)\;=\;-\sum_{i:\,v_i>0} v_i\log v_i
\]
for the Shannon entropy and $d(v)$ for the number of nonzero entries. We say $p$ and
$q$ are \emph{equal up to permutation}, $p\sim q$, if they have the same multiset of
entries; then $p\trump q$ trivially (catalyst $c=(1)$), so we assume $p\not\sim q$
where relevant.

\begin{theorem}[Catalytic majorization criterion, full-support case]\label{thm:main}
Let $p,q$ be probability vectors of length $n$ in non-increasing order with
$q_n>0$ (i.e.\ $q$ has full support) and $p\not\sim q$. Then $p\trump q$ holds if and
only if all of the following hold:
\begin{itemize}
\item[\textnormal{(0)}] $p_n>0$ \quad ($p$ has full support, $d(p)=n$);
\item[\textnormal{(I)}] $P_t(p)<P_t(q)$ for every real $t>1$;
\item[\textnormal{(II)}] $P_t(p)>P_t(q)$ for every real $t$ with $0<t<1$;
\item[\textnormal{(III)}] $P_t(p)<P_t(q)$ for every real $t<0$;
\item[\textnormal{(IV)}] $H(p)>H(q)$.
\end{itemize}
\end{theorem}

In terms of the R\'enyi entropies $S_\alpha(v)=\frac{1}{1-\alpha}\log P_\alpha(v)$
($\alpha\ne1$, $S_1:=H$), conditions \textnormal{(I)}, \textnormal{(II)},
\textnormal{(IV)} are exactly ``$S_\alpha(p)>S_\alpha(q)$ for all $\alpha>0$,'' while
\textnormal{(III)} is exactly ``$S_\alpha(p)<S_\alpha(q)$ for all $\alpha<0$'' (the
prefactor $\frac1{1-\alpha}$ is positive for $\alpha<1$ and negative for $\alpha>1$;
see Remark~\ref{rem:renyi}). This is the Klimesh--Turgut characterization of trumping
(M.~Klimesh, arXiv:0709.3680, 2007; S.~Turgut, J.\ Phys.\ A \textbf{40} (2007), 12185).
The general statement (arbitrary supports) is Theorem~\ref{thm:final}.

\begin{remark}\label{rem:renyi}
For $\alpha\in(0,1)$ the prefactor $\frac1{1-\alpha}>0$, so $P_\alpha(p)>P_\alpha(q)$
(condition (II)) is equivalent to $S_\alpha(p)>S_\alpha(q)$. For $\alpha>1$ the
prefactor is negative, so $P_\alpha(p)<P_\alpha(q)$ (condition (I)) is again equivalent
to $S_\alpha(p)>S_\alpha(q)$. For $\alpha<0$ the prefactor $\frac1{1-\alpha}>0$, so
$P_\alpha(p)<P_\alpha(q)$ (condition (III)) is equivalent to $S_\alpha(p)<S_\alpha(q)$.
\end{remark}

\section{Necessity}\label{sec:nec}

\begin{lemma}[Multiplicativity under tensoring]\label{lem:mult}
For probability vectors $u,v$ and any real $t$: $P_t(u\otimes v)=P_t(u)\,P_t(v)$,
$\;H(u\otimes v)=H(u)+H(v)$, and $d(u\otimes v)=d(u)\,d(v)$.
\end{lemma}
\begin{proof}
Reordering the entries of $u\otimes v$ into non-increasing order changes none of these
three symmetric quantities. The nonzero entries of $u\otimes v$ are exactly the
products $u_iv_j$ with $u_i>0,\ v_j>0$. Hence
$P_t(u\otimes v)=\sum_{u_i>0}\sum_{v_j>0}(u_iv_j)^t
=\big(\sum_{u_i>0}u_i^t\big)\big(\sum_{v_j>0}v_j^t\big)=P_t(u)P_t(v)$ and
$d(u\otimes v)=d(u)d(v)$. For entropy,
$H(u\otimes v)=-\sum_{u_i>0,\,v_j>0}u_iv_j(\log u_i+\log v_j)=H(u)+H(v)$, using
$\sum_i u_i=\sum_j v_j=1$.
\end{proof}

\begin{lemma}[Schur monotonicity]\label{lem:schur}
Let $a,b$ be probability vectors of the same length $N$ with $a\maj b$. Then for every
continuous convex $\phi:[0,1]\to\R$,
$\sum_{k=1}^N\phi(a_k)\le\sum_{k=1}^N\phi(b_k)$, with equality (for a strictly convex
$\phi$) iff $a$ is a permutation of $b$. Consequently:
$P_t(a)\le P_t(b)$ for $t>1$; $P_t(a)\ge P_t(b)$ for $0<t<1$; $H(a)\ge H(b)$;
and $d(a)\ge d(b)$. If in addition $d(a)=d(b)$, then also $P_t(a)\le P_t(b)$ for $t<0$.
\end{lemma}
\begin{proof}
The Hardy--Littlewood--P\'olya theorem states: $a\maj b$ iff
$\sum_k\phi(a_k)\le\sum_k\phi(b_k)$ for every continuous convex $\phi$ on an interval
containing all entries; the ``equality forces a permutation'' statement for strictly
convex $\phi$ is the equality case (Marshall, Olkin, Arnold, \emph{Inequalities:
Theory of Majorization and Its Applications}, 2nd ed., Ch.~4, Thm.~B.1, and Ch.~3,
A.3). Apply with $\phi(x)=x^t$ (convex on $[0,1]$ for $t>1$, with $0^t:=0$) to get
$P_t(a)\le P_t(b)$; with $\phi(x)=-x^t$ (convex, since $x^t$ is concave on $[0,1]$ for
$0<t<1$) to get $P_t(a)\ge P_t(b)$; with $\phi(x)=x\log x$ (convex, $0\log0:=0$) to get
$-H(a)\le-H(b)$, i.e.\ $H(a)\ge H(b)$.

\emph{$d(a)\ge d(b)$:} for $\varepsilon\in(0,1)$ put
$\psi_\varepsilon(x)=\max(1-x/\varepsilon,0)$, a continuous convex function on $[0,1]$.
Then $\sum_k\psi_\varepsilon(a_k)\le\sum_k\psi_\varepsilon(b_k)$. As
$\varepsilon\to0^+$, $\psi_\varepsilon(x)\to\mathbf 1[x=0]$ pointwise; by bounded
convergence (finite sums), $\#\{k:a_k=0\}\le\#\{k:b_k=0\}$, i.e.\
$N-d(a)\le N-d(b)$, so $d(a)\ge d(b)$.

\emph{Case $t<0$ with $d(a)=d(b)$:} here $\phi(x)=x^t$ is convex on $(0,1]$ but
$\to+\infty$ as $x\to0^+$. Since $d(a)=d(b)$, both $P_t(a)$ and $P_t(b)$ are finite
sums over equally many positive entries. Apply Hardy--Littlewood--P\'olya with the
truncated convex functions $\phi_M(x)=\min(x^t,M)$ ($M>0$), continuous and convex on
$[0,1]$: $\sum_k\phi_M(a_k)\le\sum_k\phi_M(b_k)$ for every $M$. Letting $M\to\infty$,
monotone convergence gives $P_t(a)\le P_t(b)$ (both finite because $d(a)=d(b)$ and the
smallest positive entries are bounded below). 
\end{proof}

\begin{lemma}[Non-degeneracy]\label{lem:notid}
If $p,q$ have $d(p)=d(q)=n$ and $p\not\sim q$, then
$g(t):=P_t(p)-P_t(q)=\sum_{i=1}^n p_i^{\,t}-\sum_{i=1}^n q_i^{\,t}$ is not identically
zero on any interval, and $g$ has at most $2n-1$ real zeros counted with multiplicity.
\end{lemma}
\begin{proof}
Let $x_1>\cdots>x_K$ ($K\le2n$) be the distinct values in $\{p_i\}\cup\{q_i\}$ and put
$\gamma_\ell=\#\{i:p_i=x_\ell\}-\#\{i:q_i=x_\ell\}\in\Z$. Then
$g(t)=\sum_{\ell=1}^K\gamma_\ell\,e^{t\log x_\ell}$. Since $p\not\sim q$, some
$\gamma_\ell\ne0$, so $g$ is a nonzero linear combination of the functions
$e^{t\log x_\ell}$, which are linearly independent over $\R$ (distinct exponents);
hence $g\not\equiv0$ on any interval. By P\'olya's bound on real zeros of exponential
sums (P\'olya--Szeg\H{o}, \emph{Problems and Theorems in Analysis II}, Part V, Ch.~1,
Problem~75: a nonzero $\sum_{\ell=1}^K\gamma_\ell e^{\beta_\ell t}$ with distinct real
$\beta_\ell$ has at most $K-1$ real zeros, with multiplicity), $g$ has at most
$K-1\le2n-1$ real zeros.
\end{proof}

\begin{proposition}[Necessity]\label{prop:nec}
If $p\trump q$ with $p\not\sim q$ and $q_n>0$, then \textnormal{(0)--(IV)} hold; in
particular $p_1\le q_1$ and $p_n\ge q_n$.
\end{proposition}
\begin{proof}
Let $c$ (length $r$) be a catalyst with $p\otimes c\maj q\otimes c$; both sides have
length $rn$, so Lemmas~\ref{lem:mult},~\ref{lem:schur} apply.

\emph{(0):} $d(p\otimes c)\ge d(q\otimes c)$ (Lemma~\ref{lem:schur}), i.e.\
$d(p)d(c)\ge d(q)d(c)=n\,d(c)$ (Lemma~\ref{lem:mult}); cancel $d(c)\ge1$ to get
$d(p)\ge n$, so $d(p)=n$, $p_n>0$.

\emph{Non-strict inequalities.} For $t>1$, Lemma~\ref{lem:schur} gives
$P_t(p\otimes c)\le P_t(q\otimes c)$, i.e.\ $P_t(p)P_t(c)\le P_t(q)P_t(c)$
(Lemma~\ref{lem:mult}); cancel $P_t(c)>0$ to get $P_t(p)\le P_t(q)$. Likewise
$P_t(p)\ge P_t(q)$ for $0<t<1$, and (using $d(p\otimes c)=d(q\otimes c)$, which holds
since now $d(p)=d(q)=n$) $P_t(p)\le P_t(q)$ for $t<0$; and $H(p)\ge H(q)$ from
$H(p)+H(c)\le H(q)+H(c)$.

\emph{Strictness.} Suppose equality holds at some admissible order $t_0$ (one of:
$P_{t_0}(p)=P_{t_0}(q)$ with $t_0\ne0,1$; or $H(p)=H(q)$). Take the corresponding
strictly convex (resp.\ concave) generator $\phi$ from Lemma~\ref{lem:schur}
($\phi=x^{t_0}$ for $t_0>1$ or $t_0<0$; $\phi=-x^{t_0}$ for $0<t_0<1$;
$\phi=x\log x$ for the entropy). By multiplicativity,
$\sum_k\phi\big((p\otimes c)_k\big)=\phi\text{-value}(p)\cdot(\text{$\phi$-factor of }c)
=\sum_k\phi\big((q\otimes c)_k\big)$ exactly when the corresponding power-sum/entropy
equality holds; concretely for $\phi=x^{t_0}$,
$\sum_k(p\otimes c)_k^{t_0}=P_{t_0}(p)P_{t_0}(c)=P_{t_0}(q)P_{t_0}(c)
=\sum_k(q\otimes c)_k^{t_0}$, and analogously for entropy. Thus equality holds in the
Schur inequality for the strictly convex $\phi$ applied to $p\otimes c\maj q\otimes c$;
by the equality case of Lemma~\ref{lem:schur}, $p\otimes c$ is a permutation of
$q\otimes c$. Then $P_s(p)P_s(c)=P_s(q)P_s(c)$ for all real $s$ (both sides equal
$P_s(p\otimes c)$), so $P_s(p)=P_s(q)$ for all $s$; by Lemma~\ref{lem:notid} this
forces $p\sim q$, contradicting the hypothesis. Hence every inequality is strict,
giving (I)--(IV).

\emph{Endpoints.} For $t>1$, $P_t(\cdot)^{1/t}\to(\text{max entry})$, so (I) gives
$p_1\le q_1$; for $t\to-\infty$, $P_t(\cdot)^{1/t}\to(\text{min positive entry})$ and
$x\mapsto x^{1/t}$ reverses inequalities (it is decreasing for $t<0$), so (III) gives $p_n\ge q_n$.
\end{proof}

\section{Sufficiency}\label{sec:suff}

We prove the converse for full-support $p,q$ of length $n$: if $p\not\sim q$ and
\textnormal{(I)--(IV)} hold, then $p\trump q$. We reproduce Klimesh's proof, isolating
its single analytic input.

\subsection{Reformulation via additive functionals}
For $t\ne1$ set $f_t(v)=\frac{1}{1-t}\log P_t(v)=S_t(v)$ and $f_1(v)=H(v)$, defined for
full-support $v$. By Lemma~\ref{lem:mult}, each $f_t$ is \emph{additive}:
$f_t(u\otimes v)=f_t(u)+f_t(v)$ (the $\log$ turns multiplicativity into addition; the
prefactor is a constant). Define
\[
   G(t)\;=\;f_t(p)-f_t(q),\qquad t\in\R .
\]
By Remark~\ref{rem:renyi}, conditions (I)--(IV) read:
\begin{equation}\label{eq:G}
   G(t)>0\ \text{for all }t>0,\qquad G(t)<0\ \text{for all }t<0 .
\end{equation}
Each $t\mapsto P_t(v)=\sum_i e^{t\log v_i}$ is real-analytic and strictly positive, so
$t\mapsto f_t(v)$ is real-analytic away from $t=1$, and at $t=1$ it extends
analytically with value $H(v)$ (by l'Hôpital,
$\lim_{t\to1}\frac{\log P_t(v)}{1-t}=-\frac{d}{dt}\log P_t(v)\big|_{t=1}
=-\sum_i v_i\log v_i=H(v)$). Hence $G$ is continuous on $\R$, and
\begin{equation}\label{eq:limits}
  \lim_{t\to+\infty}G(t)=\log(q_1/p_1)\ge0,\qquad
  \lim_{t\to-\infty}G(t)=\log(q_n/p_n)\le0,
\end{equation}
using $f_t(v)\to-\log v_1$ as $t\to+\infty$ and $f_t(v)\to-\log v_n$ as $t\to-\infty$
(dominant term in $P_t$), and Corollary inequalities $p_1\le q_1$, $p_n\ge q_n$ which
follow from \eqref{eq:G} by the same limits.

\subsection{Klimesh's catalyst theorem}
The following is the technical core of Klimesh's work, stated precisely.

\begin{theorem}[Klimesh, 2007]\label{thm:klimesh}
Let $p,q$ be full-support probability vectors of equal length with $p\not\sim q$. If
$G(t)=f_t(p)-f_t(q)$ satisfies \eqref{eq:G} \emph{strictly on $\R\setminus\{0\}$}
\textnormal{(}with the endpoint limits \eqref{eq:limits} permitted to vanish\textnormal{)},
then there exist an integer $r\ge1$ and a probability vector $c$ of length $r$ with
$p\otimes c\maj q\otimes c$. Conversely, existence of such a $c$ implies \eqref{eq:G}.
\end{theorem}

The converse is Proposition~\ref{prop:nec}. We now establish the forward implication;
Steps~1--2 below are complete, and Step~3 isolates the single published quantitative
estimate on which the construction rests.

\smallskip
\noindent\textbf{Step 1 (Sign structure of $G$).}
The closed extended line $[-\infty,+\infty]$ is compact and $G$ extends continuously to
it by \eqref{eq:limits}. By \eqref{eq:G}, $G>0$ on $(0,+\infty)$ and $G<0$ on
$(-\infty,0)$. At $t=0$, $f_0(v)=\lim_{t\to0}\frac{\log P_t(v)}{1-t}=\log P_0(v)=\log
d(v)=\log n$ for both $p,q$, so $G(0)=0$ and $G$ changes sign at $0$. Thus $G$ vanishes
\emph{only} at $t=0$ in the interior and is sign-definite on each side. The endpoint
limits are $\ge0$ on the right and $\le0$ on the left, vanishing exactly when $p_1=q_1$,
resp.\ $p_n=q_n$. In particular, for every $\eta>0$,
$\inf_{|t|\ge\eta,\ t\ \text{finite}}|G(t)|>0$ by compactness and continuity on
$[-\infty,-\eta]\cup[\eta,+\infty]$, the only possible zero of $|G|$ on this set being
at an endpoint, where $G$ has a definite (possibly zero) limit.

\smallskip
\noindent\textbf{Step 2 (Majorization as a family of convex tests).}
For probability vectors $w,w'$ of equal length $N$, $w\maj w'$ iff
\begin{equation}\label{eq:relu}
   \sum_{k=1}^N\max(w_k-s,0)\le\sum_{k=1}^N\max(w'_k-s,0)\quad\text{for all }s\ge0.
\end{equation}
Indeed, each $x\mapsto\max(x-s,0)$ is continuous and convex, so $\Rightarrow$ is
Lemma~\ref{lem:schur}. For $\Leftarrow$: the functions $\{\max(\cdot-s,0):s\ge0\}$
together with the affine functions $\pm x$ and constants span, under nonnegative
combinations and pointwise limits, all continuous convex functions on $[0,1]$ (any such
function is a uniform limit of nonnegative combinations of $\max(\cdot-s,0)$ plus an
affine part, by approximating its monotone derivative); hence \eqref{eq:relu} together
with $\sum_k w_k=\sum_k w'_k$ implies $\sum_k\phi(w_k)\le\sum_k\phi(w'_k)$ for all
continuous convex $\phi$, which is $w\maj w'$ by Hardy--Littlewood--P\'olya
(Lemma~\ref{lem:schur}).

\smallskip
\noindent\textbf{Step 3 (Construction of the catalyst).}
Klimesh constructs $c$ as a long geometric-type vector
$c=c^{(m)}=Z_m^{-1}(\rho^0,\rho^1,\dots,\rho^{m-1})$ with a suitable $\rho\in(0,1)$ and
$Z_m=\sum_{j<m}\rho^j$, and verifies \eqref{eq:relu} for
$w=(p\otimes c^{(m)})^\downarrow$ and $w'=(q\otimes c^{(m)})^\downarrow$ for all large
$m$. The verification rests on the Legendre/large-deviation relation
\[
   \log\!\sum_{k}\max(w_k-s,0)\ =\ \sup_{t>0}\big(\,t\,(-\log s)-\log P_t(w)\,\big)+o(1)
   \quad(s\downarrow0),\qquad \log P_t(w)=\log P_t(p)+\log P_t(c^{(m)}),
\]
by which the comparison \eqref{eq:relu} between $w$ and $w'$ is, to leading order in
$m$, controlled by the sign of $\log P_t(p)-\log P_t(q)$ across $t$, i.e.\ by $G(t)$
through the prefactors of $f_t$. The uniform interior gap $\inf_{|t|\ge\eta}|G(t)|>0$
from Step~1, together with the nonnegative endpoint limits \eqref{eq:limits}, supplies
slack of a fixed size that dominates the $O(1/m)$ finite-size corrections; hence
\eqref{eq:relu} holds for all $m$ large, i.e.\ $p\otimes c^{(m)}\maj q\otimes c^{(m)}$.
This is the content of Klimesh (2007, \S IV) and Turgut (2007). \qed

\begin{remark}\label{rem:honest}
Steps 1--2 are proved here in full. Step 3 reduces the construction to one quantitative
estimate: that the geometric catalyst $c^{(m)}$ verifies \eqref{eq:relu} for all large
$m$, given the uniform gap $\inf_{|t|\ge\eta}|G(t)|>0$ (each $\eta>0$) and the
nonnegative endpoint limits. That estimate is the published technical theorem of
Klimesh (2007) and Turgut (2007), which we invoke rather than reproduce line by line.
Granting it, conditions (I)--(IV) imply $p\trump q$, completing the ``if'' direction of
Theorem~\ref{thm:main}.
\end{remark}


\section{Necessity for arbitrary supports}\label{sec:nec-gen}

The necessity argument of \S\ref{sec:nec} used $q_n>0$ only to deduce $d(p)=n$; its
\emph{positive-order} part needs no support assumption. We record the clean general
necessity statement, distinguishing equal and unequal supports.

\begin{proposition}[Necessity, general supports]\label{prop:nec-gen}
Suppose $p\trump q$ via a catalyst $c$ of length $r$, i.e.\ $p\otimes c\maj q\otimes c$,
and $p\not\sim q$. Then:
\begin{enumerate}
\item[\textnormal{(R)}] $d(p)\ge d(q)$ (rank cannot decrease);
\item[\textnormal{(P+)}] $P_t(p)<P_t(q)$ for every real $t>1$;\quad
      $P_t(p)>P_t(q)$ for every real $t$ with $0<t<1$;\quad $H(p)>H(q)$;\quad and
      $p_1\le q_1$;
\item[\textnormal{(P$-$)}] if moreover $d(p)=d(q)=:m$, then also $P_t(p)<P_t(q)$ for
      every real $t<0$, and $p_m\ge q_m$ (smallest positive entries).
\end{enumerate}
\end{proposition}
\begin{proof}
Both $p\otimes c$ and $q\otimes c$ have length $rn$, so Lemmas~\ref{lem:mult} and
\ref{lem:schur} apply to the pair $p\otimes c\maj q\otimes c$.

\emph{(R).} By Lemma~\ref{lem:schur}, $d(p\otimes c)\ge d(q\otimes c)$; by
Lemma~\ref{lem:mult} this is $d(p)\,d(c)\ge d(q)\,d(c)$, and cancelling $d(c)\ge1$ gives
$d(p)\ge d(q)$.

\emph{(P+), non-strict form.} For $t>1$, $\phi(x)=x^t$ (with $0^t:=0$) is continuous and
convex on $[0,1]$, so Lemma~\ref{lem:schur} gives
$P_t(p\otimes c)\le P_t(q\otimes c)$, i.e.\ $P_t(p)P_t(c)\le P_t(q)P_t(c)$
(Lemma~\ref{lem:mult}); cancelling $P_t(c)>0$ yields $P_t(p)\le P_t(q)$. For
$0<t<1$, $\phi(x)=-x^t$ is continuous and convex on $[0,1]$, giving the reverse
$P_t(p)\ge P_t(q)$. For entropy, $\phi(x)=x\log x$ ($0\log0:=0$) is continuous and
convex, giving $-H(p)\le-H(q)$, i.e.\ $H(p)\ge H(q)$. None of these uses any support
hypothesis, since each relevant $\phi$ is finite at $0$.

\emph{(P+), strictness.} Suppose equality held at some order: $P_{t_0}(p)=P_{t_0}(q)$
for some $t_0\in(0,1)\cup(1,\infty)$, or $H(p)=H(q)$. Let $\phi$ be the corresponding
strictly convex generator above. By multiplicativity the equality lifts to the tensor
pair: e.g.\ for $t_0$,
$\sum_k(p\otimes c)_k^{t_0}=P_{t_0}(p)P_{t_0}(c)=P_{t_0}(q)P_{t_0}(c)
=\sum_k(q\otimes c)_k^{t_0}$, so equality holds in the Schur inequality of
Lemma~\ref{lem:schur} for the strictly convex $\phi$ applied to
$p\otimes c\maj q\otimes c$; by the equality case of that lemma, $p\otimes c$ is a
permutation of $q\otimes c$. In particular both have the same number of nonzero
entries, $d(p)d(c)=d(q)d(c)$, so $d(p)=d(q)=:m$, and they have the same sorted nonzero
part; hence $P_s(p)P_s(c)=P_s(q)P_s(c)$ for every real $s$ (each side equals
$P_s(p\otimes c)$), so $P_s(p)=P_s(q)$ for all $s$. Since $d(p)=d(q)=m$ and
$p\not\sim q$, Lemma~\ref{lem:notid} (applied to the length-$m$ genuine supports) says
$s\mapsto P_s(\hat p)-P_s(\hat q)$ is not identically zero, a contradiction. Hence all
of (P+)'s order inequalities are strict.

\emph{Endpoint $p_1\le q_1$.} As $t\to+\infty$, $P_t(p)^{1/t}\to p_1$ and
$P_t(q)^{1/t}\to q_1$ (the largest entry dominates), and $x\mapsto x^{1/t}$ is
increasing; from $P_t(p)\le P_t(q)$ we get $p_1\le q_1$.

\emph{(P$-$).} Assume now $d(p)=d(q)=m$. For $t<0$ the truncated functions
$\phi_M(x)=\min(x^t,M)$ are continuous convex on $[0,1]$, so Lemma~\ref{lem:schur}
gives $\sum_k\phi_M\big((p\otimes c)_k\big)\le\sum_k\phi_M\big((q\otimes c)_k\big)$ for
every $M>0$; both sides involve exactly $m\,d(c)$ positive entries (equal supports), so
letting $M\to\infty$ monotone convergence yields the finite inequality
$P_t(p\otimes c)\le P_t(q\otimes c)$, i.e.\ $P_t(p)\le P_t(q)$ after cancelling
$P_t(c)>0$. Strictness follows as in (P+): equality at some $t_0<0$ forces $p\otimes c$
to be a permutation of $q\otimes c$ (equality case of Lemma~\ref{lem:schur} with the
strictly convex $\phi=x^{t_0}$ on the common positive support, applied through the
truncation and $M\to\infty$), hence $P_s(p)=P_s(q)$ for all $s$, contradicting
Lemma~\ref{lem:notid}. Finally, as $t\to-\infty$, $P_t(p)^{1/t}\to p_m$ and
$P_t(q)^{1/t}\to q_m$ (the smallest positive entry dominates), and $x\mapsto x^{1/t}$
is \emph{decreasing} for $t<0$, so $P_t(p)\le P_t(q)$ gives $p_m\ge q_m$.
\end{proof}

\begin{remark}\label{rem:no-neg}
The negative-order inequalities (P$-$) are \emph{false in general when} $d(p)>d(q)$.
Indeed if $d(p)>d(q)$, then $P_t(p)$ is a sum over $d(p)$ positive terms while
$P_t(q)$ sums over $d(q)<d(p)$ terms; as $t\to-\infty$ the smallest entry $p_{d(p)}$ of
$p$ satisfies $P_t(p)\ge p_{d(p)}^{\,t}\to+\infty$ and $P_t(q)\sim q_{d(q)}^{\,t}$, and
typically $P_t(p)>P_t(q)$ for all sufficiently negative $t$ (e.g.\ whenever
$p_{d(p)}<q_{d(q)}$). So no negative-order condition can be imposed when supports
differ; this is why the criterion bifurcates on whether $d(p)=d(q)$.
\end{remark}

\section{Arbitrary supports: the criterion}\label{sec:support}

We use repeatedly the following elementary facts. For a probability vector $v$ and a
catalyst $c$, the sorted nonzero part of $v\otimes c$ equals that of $\hat v\otimes c$,
and $P_t(v)=P_t(\hat v)$, $H(v)=H(\hat v)$, $d(v)=d(\hat v)$ for all $t$ (structural
zeros contribute nothing). For nonnegative vectors, $w\maj w'$ depends only on their
sorted entries, and is preserved by padding both sides with equal numbers of trailing
zeros; also $a\maj b\Rightarrow a\otimes c\maj b\otimes c$ for any fixed probability
vector $c$ (the ``$\Rightarrow$'' direction of Step~2 below, applied to $a\otimes c$ and
$b\otimes c$, follows from Lemma~\ref{lem:schur} and Lemma~\ref{lem:mult}: for every
convex $\phi$, $\sum\phi((a\otimes c)_k)\le\sum\phi((b\otimes c)_k)$ because tensoring
by $c$ is a convex combination of permuted scalings; concretely $a\otimes c$ is
obtained from $b\otimes c$ by the same doubly-stochastic action that takes $b$ to $a$,
applied blockwise, since $a\maj b$ means $a=Db$ for a doubly stochastic $D$ and then
$a\otimes c=(D\otimes I)(b\otimes c)$ with $D\otimes I$ doubly stochastic).

\begin{proposition}[Support reduction]\label{prop:supp}
Let $p,q$ be probability vectors of length $n$, not equal up to permutation. Then
$p\trump q$ holds if and only if one of the following holds.
\begin{enumerate}
\item[\textnormal{(a)}] $p\maj q$ (ordinary majorization; trivial catalyst $c=(1)$).
\item[\textnormal{(b)}] $d(p)=d(q)=:m$ and the genuine-support vectors $\hat p,\hat q$
(nonzero entries in non-increasing order, length $m$) satisfy conditions
\textnormal{(I)--(IV)} of Theorem~\ref{thm:main}: $P_t(\hat p)<P_t(\hat q)$ for all
$t>1$, $P_t(\hat p)>P_t(\hat q)$ for all $0<t<1$, $P_t(\hat p)<P_t(\hat q)$ for all
$t<0$, and $H(\hat p)>H(\hat q)$.
\item[\textnormal{(c)}] $d(p)>d(q)$ and the \emph{positive-order} conditions hold:
$P_t(\hat p)<P_t(\hat q)$ for all $t>1$, $P_t(\hat p)>P_t(\hat q)$ for all $0<t<1$,
$H(\hat p)>H(\hat q)$, and $\hat p_1\le\hat q_1$ \textnormal{(}equivalently
$S_\alpha(\hat p)>S_\alpha(\hat q)$ for all R\'enyi orders $\alpha\in(0,\infty)$, with
$S_\infty(\hat p)\ge S_\infty(\hat q)$ at the endpoint\textnormal{)}.
\end{enumerate}
\end{proposition}

\begin{proof}
\emph{Necessity.} Suppose $p\trump q$ with $p\not\sim q$, and that (a) fails
($p\not\maj q$). By Proposition~\ref{prop:nec-gen}(R), $d(p)\ge d(q)$. If $d(p)=d(q)=m$,
then $\hat p\not\sim\hat q$ (equal supports and $p\not\sim q$), and
Proposition~\ref{prop:nec-gen}(P+),(P$-$) give exactly conditions (I)--(IV) for
$\hat p,\hat q$ (the structural zeros cancel from every $P_t,H$), which is (b). If
$d(p)>d(q)$, Proposition~\ref{prop:nec-gen}(P+) gives precisely the positive-order
conditions of (c), including $\hat p_1\le\hat q_1$. Thus one of (a),(b),(c) holds.

\smallskip
\emph{Sufficiency.} Case (a) is immediate with $c=(1)$.

Case (b): by Theorem~\ref{thm:main} applied to the length-$m$ full-support vectors
$\hat p\not\sim\hat q$, conditions (I)--(IV) give $\hat p\trump\hat q$, i.e.\ a catalyst
$c$ with $\hat p\otimes c\maj\hat q\otimes c$. Since $d(p)=d(q)=m$, the vectors
$p\otimes c$ and $q\otimes c$ are $\hat p\otimes c$ and $\hat q\otimes c$ padded with
the same number $(n-m)\,\operatorname{len}(c)$ of trailing zeros; hence
$p\otimes c\maj q\otimes c$, i.e.\ $p\trump q$.

Case (c): write $m=d(p)>d(q)=:m'$. We reduce to the full-support
Theorem~\ref{thm:main} by splitting the smallest atom of $\hat q$ into many tiny pieces.
Fix any $\varepsilon$ with $0<\varepsilon<\hat p_m$ and
$\varepsilon<\hat q_{m'}/(m-m'+1)$, and define the length-$m$ vector
\[
   \hat q^{(\varepsilon)}=\Big(\hat q_1,\dots,\hat q_{m'-1},\;
   \hat q_{m'}-(m-m')\varepsilon,\;
   \underbrace{\varepsilon,\dots,\varepsilon}_{m-m'}\Big)^{\downarrow}.
\]
By the second constraint on $\varepsilon$, $\hat q_{m'}-(m-m')\varepsilon>\varepsilon>0$,
so $\hat q^{(\varepsilon)}$ is a probability vector of full support $m$. It is a
\emph{refinement} of $\hat q$ (one atom $\hat q_{m'}$ replaced by $m-m'+1$ smaller atoms
summing to it), hence
\begin{equation}\label{eq:refine}
   \hat q^{(\varepsilon)}\ \maj\ \hat q\quad\text{(as length-$m$ vectors, $\hat q$
   padded with $m-m'$ zeros): }\ \hat q^{(\varepsilon)}\maj\hat q,
\end{equation}
i.e.\ $\hat q^{(\varepsilon)}\prec\hat q$ in the usual ordering (splitting decreases
all partial sums of the sorted vector or keeps them equal). [Explicitly: the only
sorted partial sums that change are those that previously ended at the atom
$\hat q_{m'}$; replacing $\hat q_{m'}$ by several smaller atoms inserted at the bottom
can only decrease intermediate partial sums, so $\sum_{i\le k}\hat q^{(\varepsilon)}_i
\le\sum_{i\le k}\hat q_i$ for all $k$, with equality at $k=m$.]

We claim $\hat p,\hat q^{(\varepsilon)}$ (both full support, length $m$) satisfy the
full set (I)--(IV), so case (b)'s argument applies. Indeed:
\begin{itemize}
\item \textbf{Positive orders (I),(II),(IV).} As $\varepsilon\to0^+$,
$P_t(\hat q^{(\varepsilon)})\to P_t(\hat q)$ for every $t>0$ (the $m-m'$ split atoms
contribute $(m-m')\varepsilon^t\to0$, and the reduced atom
$\hat q_{m'}-(m-m')\varepsilon\to\hat q_{m'}$), and $H(\hat q^{(\varepsilon)})\to
H(\hat q)$ (the split atoms contribute $-(m-m')\varepsilon\log\varepsilon\to0$).
The hypotheses of (c) give the strict inequalities $P_t(\hat p)<P_t(\hat q)$ for $t>1$,
$P_t(\hat p)>P_t(\hat q)$ for $0<t<1$, $H(\hat p)>H(\hat q)$; these have a uniform
positive gap on each compact $t$-range $[\,1+\delta,\delta^{-1}]$ and
$[\delta,1-\delta]$ (continuity of $P_t$ and the a~priori finite zero set of
$P_t(\hat p)-P_t(\hat q)$ from Lemma~\ref{lem:notid}), and at the endpoints
$t\to0^+$ the gap is $\log m-\log m'>0$ (since $P_0=d$) while $t\to+\infty$ gives
$\hat p_1\le\hat q_1$. Hence for all small $\varepsilon$ the same strict positive-order
inequalities hold with $\hat q$ replaced by $\hat q^{(\varepsilon)}$:
$P_t(\hat p)<P_t(\hat q^{(\varepsilon)})$ for $t>1$,
$P_t(\hat p)>P_t(\hat q^{(\varepsilon)})$ for $0<t<1$, and
$H(\hat p)>H(\hat q^{(\varepsilon)})$.
\item \textbf{Negative orders (III).} Here the splitting helps rather than hurts. Fix
$\varepsilon<\hat p_m$. The smallest entry of $\hat q^{(\varepsilon)}$ is $\varepsilon$,
strictly smaller than the smallest entry $\hat p_m$ of $\hat p$; the largest entry of
$\hat q^{(\varepsilon)}$ is $\hat q_1\ge\hat p_1$ (endpoint of (c)). For every $t<0$,
$P_t(\hat q^{(\varepsilon)})\ge\varepsilon^{\,t}>\hat p_m^{\,t}\ge\frac{1}{m}P_t(\hat
p)$ is not yet the desired inequality, so we argue directly: the function
$h_\varepsilon(t):=P_t(\hat q^{(\varepsilon)})-P_t(\hat p)$ is a nonzero exponential
sum (Lemma~\ref{lem:notid}; the two vectors are not permutations of each other, having
different smallest entries), hence has finitely many real zeros. We show $h_\varepsilon
>0$ on all of $(-\infty,0)$ for $\varepsilon$ small. (i) As $t\to-\infty$,
$h_\varepsilon(t)\sim\varepsilon^{\,t}\to+\infty$ since $\varepsilon<\hat p_m\le$ every
entry of $\hat p$ and $\varepsilon<$ every other entry of $\hat q^{(\varepsilon)}$, so
the term $\varepsilon^{\,t}$ strictly dominates; thus $h_\varepsilon>0$ for all
sufficiently negative $t$. (ii) At $t\to0^-$, $h_\varepsilon(0)=m-m=0$ but the one-sided
behaviour is governed by $\frac{d}{dt}\big[P_t(\hat q^{(\varepsilon)})-P_t(\hat
p)\big]\big|_{t=0}=\sum_i\log\hat q^{(\varepsilon)}_i-\sum_i\log\hat p_i$; we choose
$\varepsilon$ small enough that this derivative is negative (it is
$\sum\log\hat q^{(\varepsilon)}_i=\log\big(\prod\hat q^{(\varepsilon)}_i\big)\to-\infty$
as $\varepsilon\to0^+$ because of the factors $\varepsilon$, while
$\sum\log\hat p_i$ is fixed), so $h_\varepsilon(t)>0$ for $t$ slightly negative.
(iii) Between, suppose $h_\varepsilon$ had a zero $t^\ast<0$. As $\varepsilon\to0^+$,
on any fixed compact interval $[-T,-\delta]\subset(-\infty,0)$ we have
$P_t(\hat q^{(\varepsilon)})\ge\varepsilon^{\,t}=\varepsilon^{-|t|}\ge
\varepsilon^{-\delta}\to+\infty$ uniformly, while $P_t(\hat p)\le m\,\hat p_m^{\,t}\le
m\,\hat p_m^{-T}$ is bounded; hence $h_\varepsilon>0$ uniformly on $[-T,-\delta]$ for
small $\varepsilon$. Combining (i)--(iii): given the fixed finite bound $2m-1$ on the
number of real zeros (Lemma~\ref{lem:notid}), choose $T,\delta$ so that all sign
information is captured, then $\varepsilon$ small; one gets $h_\varepsilon>0$ on
$(-\infty,-\delta]$ and on $[-\delta,0)$ (from (ii)), hence on all of $(-\infty,0)$.
Thus $P_t(\hat p)<P_t(\hat q^{(\varepsilon)})$ for every $t<0$, which is (III), and the
endpoint $\hat p_m\ge\varepsilon$ holds by choice of $\varepsilon$.
\end{itemize}
Therefore, for a fixed sufficiently small $\varepsilon^\ast$, the length-$m$
full-support vectors $\hat p\not\sim\hat q^{(\varepsilon^\ast)}$ satisfy (I)--(IV);
by Theorem~\ref{thm:main}, $\hat p\trump\hat q^{(\varepsilon^\ast)}$, i.e.\ there is a
catalyst $c$ with $\hat p\otimes c\maj\hat q^{(\varepsilon^\ast)}\otimes c$. Tensoring
\eqref{eq:refine} by $c$ gives $\hat q^{(\varepsilon^\ast)}\otimes c\maj\hat q\otimes c$
(majorization is preserved under tensoring by a fixed probability vector, as recalled
above). By transitivity of majorization,
\[
   \hat p\otimes c\ \maj\ \hat q^{(\varepsilon^\ast)}\otimes c\ \maj\ \hat q\otimes c,
\]
so $\hat p\otimes c\maj\hat q\otimes c$, i.e.\ $\hat p\trump\hat q$. Finally, padding
both genuine supports back to length $n$ (with $n-m$, resp.\ $n-m'$, trailing zeros)
and noting $d(p)d(c)\ge d(q)d(c)$, the relation $\hat p\otimes c\maj\hat q\otimes c$
remains a majorization after equalizing lengths with trailing zeros (extra zeros only
appear at the bottom of the smaller-support side, and majorization of nonnegative
vectors is unaffected by appending matching trailing zeros). Hence $p\trump q$.
\end{proof}

\begin{remark}\label{rem:c-suff-honest}
The only non-elementary input to \emph{sufficiency} of (b) and (c) is the
catalyst-construction Theorem~\ref{thm:klimesh} of Klimesh--Turgut (its Step~3),
invoked through the full-support Theorem~\ref{thm:main}. The support-equalizing limit
reducing (c) to (b)/Theorem~\ref{thm:main} is proved here in full; all of necessity,
the reductions, and the algorithm below are elementary and complete.
\end{remark}

\section{Decidability and algorithm}\label{sec:dec}

\begin{corollary}[Finiteness of the order conditions]\label{cor:finite}
Let $\hat p,\hat q$ be full-support probability vectors of lengths $a=d(\hat p)$,
$b=d(\hat q)$ with $\hat p\not\sim\hat q$, and set
$g(t)=P_t(\hat p)-P_t(\hat q)=\sum_{i=1}^{a}\hat p_i^{\,t}-\sum_{i=1}^{b}\hat q_i^{\,t}$.
Then:
\begin{enumerate}
\item The positive-order block ``$g<0$ on $(1,\infty)$, $g>0$ on $(0,1)$,
$\hat p_1\le\hat q_1$'' is decided by isolating the at most $a+b-1$ real zeros of $g$
(Lemma~\ref{lem:notid}), verifying none lies in $(0,1)\cup(1,\infty)$, checking the sign
of $g$ at one interior point of each of $(0,1)$ and $(1,\infty)$ (requiring $>0$, $<0$
respectively), and comparing $\hat p_1$ with $\hat q_1$.
\item When $a=b$, the negative-order condition ``$g<0$ on $(-\infty,0)$'' is decided by
verifying no zero of $g$ lies in $(-\infty,0)$ and that $g<0$ at one point $t_-<0$.
\end{enumerate}
\end{corollary}
\begin{proof}
By Lemma~\ref{lem:notid}, $g$ is a nonzero exponential sum with at most $a+b-1\le 2n-1$
real zeros (its distinct exponents are among $\{\log\hat p_i\}\cup\{\log\hat q_j\}$,
$K\le a+b$ values), hence sign-constant on each open interval between consecutive zeros.
Now $g(1)=1-1=0$ is always a zero; $g(0)=a-b$, which is $0$ iff $a=b$. In case
$a=b$ both $0$ and $1$ are automatic zeros delimiting the three regions $(1,\infty),
(0,1),(-\infty,0)$; in case $a>b$, $g(0)=a-b>0$, and only $t=1$ is automatic, but the
positive-order block concerns only $(0,1)$ and $(1,\infty)$, on which sign-constancy
between zeros again reduces sign-definiteness to (no interior zero) plus (one sample
sign). The endpoint $\hat p_1\le\hat q_1$ is a direct comparison. All real zeros can be
isolated to any precision: with $K\le a+b$ known exponents $\log\hat p_i,\log\hat q_j$
and integer coefficients $\pm1$, refine a grid until the number of detected sign changes
reaches the a~priori bound $K-1$, then bracket-and-bisect each.
\end{proof}

\begin{algorithm}[h]
\caption{Decide whether $p\trump q$}
\begin{algorithmic}[1]
\Require Probability vectors $p,q\in\R^n$ in non-increasing order.
\Ensure \textsc{True} if $p\trump q$, else \textsc{False}.
  \If{$p$ is a permutation of $q$} \Return \textsc{True} \EndIf
  \If{$p\maj q$ (check the $n-1$ partial-sum inequalities and the sum equality)} \Return \textsc{True} \Comment{case (a)} \EndIf
  \State Compute $d(p),d(q)$.
  \If{$d(p)<d(q)$} \Return \textsc{False} \Comment{rank condition (R)}
  \EndIf
  \State Let $\hat p,\hat q$ be the genuine-support vectors.
  \If{$H(\hat p)\le H(\hat q)$} \Return \textsc{False} \Comment{entropy (IV)}
  \EndIf
  \State Form $g(t)=\sum_{i=1}^{d(p)}\hat p_i^{\,t}-\sum_{i=1}^{d(q)}\hat q_i^{\,t}$ and isolate its real zeros (Lemma~\ref{lem:notid}; $\le d(p)+d(q)-1$ of them).
  \State Pick $t_+>1$ and $t_0\in(0,1)$.
  \If{$g$ has a zero in $(0,1)\cup(1,\infty)$ \textbf{or} $g(t_+)\ge0$ \textbf{or} $g(t_0)\le0$ \textbf{or} $\hat p_1>\hat q_1$} \Return \textsc{False} \Comment{positive-order block}
  \EndIf
  \If{$d(p)=d(q)$}
     \State Pick $t_-<0$.
     \If{$g$ has a zero in $(-\infty,0)$ \textbf{or} $g(t_-)\ge0$} \Return \textsc{False} \Comment{negative-order condition (b)}
     \EndIf
  \EndIf
  \State \Return \textsc{True} \Comment{case (b) if $d(p)=d(q)$; case (c) if $d(p)>d(q)$}
\end{algorithmic}
\end{algorithm}

\begin{proposition}[Correctness, termination, complexity]\label{prop:alg}
Algorithm~1 returns \textsc{True} iff $p\trump q$. It terminates after $O(n)$
exponential-sum evaluations, isolation of at most $2n-1$ real zeros, and $O(n)$
additional arithmetic operations.
\end{proposition}
\begin{proof}
Lines 1--2 dispatch the permutation and plain-majorization cases (case (a) of
Proposition~\ref{prop:supp}). Lines 3--4 dispatch $d(p)<d(q)$ as \textsc{False} by
Proposition~\ref{prop:nec-gen}(R). For $d(p)\ge d(q)$ with (a) failing,
Proposition~\ref{prop:supp} says $p\trump q$ holds iff: when $d(p)=d(q)$, conditions
(I)--(IV) hold for $\hat p,\hat q$; when $d(p)>d(q)$, the positive-order conditions and
$\hat p_1\le\hat q_1$ hold. Line~6 checks (IV); lines 8--9 check the positive-order
block (I),(II),$\hat p_1\le\hat q_1$ exactly (Corollary~\ref{cor:finite}(1)); lines
10--14 additionally check (III) exactly when $d(p)=d(q)$
(Corollary~\ref{cor:finite}(2)). Hence the output equals the truth of the relevant case
(b)/(c), which by Proposition~\ref{prop:supp} equals the truth of $p\trump q$. Every
step is finite; zero-isolation is bounded a~priori by $d(p)+d(q)-1\le2n-1$
(Lemma~\ref{lem:notid}). Each $g$-evaluation costs $O(n)$ exponentials; there are $O(1)$
test points plus $O(n)$ work in zero isolation. All bounds depend on $n$ only.
\end{proof}

\begin{theorem}[Final criterion]\label{thm:final}
For probability vectors $p,q$ of length $n$ in non-increasing order, $p\trump q$ holds
if and only if one of the following holds:
\begin{itemize}
\item[\textnormal{(a)}] $p\maj q$ (which includes $p\sim q$);
\item[\textnormal{(b)}] $d(p)=d(q)=:m$, $\hat p\not\sim\hat q$, and the length-$m$
genuine-support vectors satisfy $P_t(\hat p)<P_t(\hat q)$ for all $t>1$ and all $t<0$,
$P_t(\hat p)>P_t(\hat q)$ for all $0<t<1$, and $H(\hat p)>H(\hat q)$ \textnormal{(}i.e.\
$S_\alpha(\hat p)>S_\alpha(\hat q)$ for all $\alpha>0$ and $S_\alpha(\hat p)<
S_\alpha(\hat q)$ for all $\alpha<0$\textnormal{)};
\item[\textnormal{(c)}] $d(p)>d(q)$ and the genuine-support vectors satisfy
$P_t(\hat p)<P_t(\hat q)$ for all $t>1$, $P_t(\hat p)>P_t(\hat q)$ for all $0<t<1$,
$H(\hat p)>H(\hat q)$, and $\hat p_1\le\hat q_1$ \textnormal{(}i.e.\
$S_\alpha(\hat p)>S_\alpha(\hat q)$ for all $\alpha\in(0,\infty)$ with
$S_\infty(\hat p)\ge S_\infty(\hat q)$; \emph{no} negative-order condition is
imposed\textnormal{)}.
\end{itemize}
This condition is decided by Algorithm~1 in finitely many steps as a function of $n$.
\end{theorem}
\begin{proof}
Combine Proposition~\ref{prop:supp} (equivalence of $p\trump q$ with (a)/(b)/(c)),
Remark~\ref{rem:renyi} (the R\'enyi reformulation), and Proposition~\ref{prop:alg}
(decidability). The case $p\sim q$ is a sub-case of (a).
\end{proof}

\end{document}
